\section{Introduction: the missing test}
\label{sec:intro}
Spatial transcriptomics (ST) shows where genes are expressed inside tissue~\cite{spatialtx,xenium}.
It is informative but expensive. H\&E slides are routine and much cheaper, so many methods now try to
predict spatial expression directly from histology. They use regression, transformers, spatial graphs,
retrieval, multi-resolution features, or generative models
~\cite{stnet,histogene,hist2st,bleep,triplex,stflow}.

The usual test holds out a patient but keeps the organ the same. A breast test slide, for example, is
predicted by a model that has already seen other breast slides. This measures performance on a new
patient, but not on new tissue. Our question is simple: \emph{what happens when the test organ was never
seen during training?} This matters because paired ST data are scarce and uneven across organs.

This test is easy to get wrong. If genes are selected using all slides, the test organ influences what
is evaluated. If genes are correctly selected from training slides only, some may be almost silent in
the new organ. Pearson correlation then has little variation to measure. A low score can reflect model
error, a mismatched gene panel, or both. Reporting only one average hides this distinction.

We introduce \coast{} (\textbf{C}ross-\textbf{O}rgan \textbf{A}ssessment of \textbf{S}patial
\textbf{T}ranscriptomics prediction). It separates new-patient prediction (POOLED) from new-organ
transfer (LOOO), selects genes from training data only, and reports detectability and target variance
beside accuracy. We apply it to HEST-1k~\cite{hest} and STImage-1K4M~\cite{stimage}. Every method uses
the same frozen UNI image features~\cite{uni}, allowing us to compare predictors rather than encoders.
Figure~\ref{fig:overview} summarizes the split, the shared intervention, and the main result.

We ask four questions:
\begin{enumerate}
  \item How much performance is lost when the test organ is unseen?
  \item Do model rankings change between new-patient and new-organ evaluation?
  \item When does Pearson correlation stop separating models reliably?
  \item Does the same morphology conditioner help different architectures?
\end{enumerate}

The fourth question makes the benchmark diagnostic, not only descriptive. In the main conditioning
study, FiLM improves six representative regression and spatial backbones in both regimes, including
a recent multi-level Gene-DML-style predictor. The
appendix reports the complete architecture screening, including cases where the intervention is neutral
or harmful. This gives a clear result: morphology conditioning often helps, but its value depends on
the prediction architecture.

Our contributions are:
\begin{itemize}
  \item \textbf{A clean cross-organ test:} patient-disjoint POOLED and LOOO splits, genes ranked on
  training data only, equal weighting of organs, and released evaluation artifacts.
  \item \textbf{A metric diagnosis:} low gene detection and target variance identify cases where
  Pearson correlation cannot meaningfully separate models.
  \item \textbf{A controlled intervention:} the same morphology conditioner improves six
  representative architectures in both regimes; complete screening is reported in the appendix.
\end{itemize}
